\subsection{The same balance can survive different source identities}

Apple's following-quarter filing provides an actual counterpart to the
constructed examples. Its balance sheet repeats the 28~September~2024
comparative balances already examined in the company case. The retained
annual HTML is a converted presentation without inline tags, so it cannot
supply the proposed inline trace. Instead, the selected comparison follows
the next-quarter inline comparative amounts to the retained annual XBRL
instance. This compares reports of the same endpoint across two filings; it
does not show that one file was extracted from the other
\citep{apple2025q1,apple2024xbrlinstance,apple2024convertedhtml}.

In the interim filing, the comparative total-assets occurrence has identifier
\code{f-135}, value text \code{364,980}, scale 6 and declared decimals $-6$.
Its format names \code{num-dot-decimal} in the 12~February~2020 transformation
registry. The selected transform removes the grouping comma and canonicalizes
the numeric string; scaling gives $364{,}980\times10^6=364{,}980{,}000{,}000$
US dollars. That reported amount agrees with annual-instance occurrence
\code{f-169}. Scale converts the displayed millions to dollars. Declared
decimals retain the source's million-dollar reporting accuracy; adding zeros
to the string does not reveal an exact unrounded asset value
\citep{apple2025q1,apple2024xbrlinstance,arelleinlinevalues2026}.

The context identifiers differ: the interim uses \code{c-18}, while the annual
instance uses \code{c-21}. Comparing the identifiers would reject a valid
candidate match. Inspecting their contents instead shows the same CIK entity
identifier and scheme, the same 28~September~2024 instant and no segment,
scenario or dimensional additions in these selected contexts. The expanded
concept name, USD unit and declared accuracy also agree. This is a checked
comparison of these simple contexts, not a general context-equivalence
algorithm or schema-validation result
\citep{apple2025q1,apple2024xbrlinstance}.

The liabilities and equity comparisons use the same conditions and give
308,030 and 56,950 million dollars. Their reported totals reconcile:
$308{,}030+56{,}950=364{,}980$. Equity, however, appears twice in the interim:
\code{f-173} is the comparative balance-sheet amount and \code{f-176} is the
following period's beginning equity. Both carry 56,950 million dollars under
the same selected concept, context and unit. The annual instance also retains
two matching occurrences, \code{f-207} and \code{f-256}. Amount and context
matching therefore identifies a group, not a unique pairing of source
locations. Adding the two interim occurrences would produce 113,900 million
dollars and would double the reported equity balance. A balance query uses
one accepted amount while preserving both locations and their statement roles
\citep[printed pp.~3--4]{apple2025q1}; the annual counterparts are retained in the instance \citep{apple2024xbrlinstance}.

The accumulated other comprehensive loss exposes a different failure. Interim
occurrence \code{f-171} contains numeric text \code{7,172}, with scale 6 and
an explicit minus-sign attribute. Its reported value is consequently
$-7{,}172{,}000{,}000$ dollars, matching annual occurrence \code{f-205}. Reading
the digits alone loses the sign; reading the surrounding parentheses and then
applying the sign attribute again would reverse it twice. The source text,
markup attribute and displayed accounting reading must remain distinguishable
\citep[printed p.~3]{apple2025q1}; the signed annual value is retained in the instance \citep{apple2024xbrlinstance}.

The bounded inventory retains all 24 interim and 32 annual-instance occurrences
of the four selected concepts, then selects five occurrences from each source
at the unsegmented September endpoint. Other dates and dimensional populations
remain visible as outside this comparison. All selected reported amounts
agree under the stated structural and arithmetic checks. The agreement supplies
an actual preservation example; it establishes neither historical availability
of the retrieved bytes nor schema validity, complete filing coverage or an
executed adapter. Those distinct questions remain part of the wider acceptance
work \citep{apple2025q1,apple2024xbrlinstance}.
